\chapter{Einen Mod veröffentlichen}
\label{ch:quality/releasing}

Wenn du einen Mod veröffentlichst, können andere ihn nutzen. Ein
veröffentlichter Mod wird womöglich in Einzelspielerwelten oder auf
Servern verwendet, auch auf öffentlichen.

\section{Eine Lizenz wählen}\label{eine-lizenz-wuxe4hlen}

Du musst für deinen Mod eine Lizenz\index{lizenz@Lizenz} angeben. Sie
sagt anderen, auf welche Weise sie dein Werk nutzen dürfen. Hat dein Mod
keine Lizenz, wissen andere nicht, ob sie ihn verändern, weitergeben
oder auf einem öffentlichen Server einsetzen dürfen.

Code und Grafiken stellen unterschiedliche Anforderungen an ihre
Lizenzen. Creative-Commons-Lizenzen solltest du zum Beispiel nicht für
Quellcode verwenden; für künstlerische Werke wie Bilder, Texte und
Meshes\index{mesh@Mesh} können sie aber gut geeignet sein.

Du darfst jede Lizenz verwenden. Mods, die keine Bearbeitungen erlauben,
sind im offiziellen Luanti-Forum allerdings nicht zugelassen. (Damit ein
Mod im Forum erlaubt ist, müssen andere Entwickler ihn verändern und die
veränderte Fassung veröffentlichen dürfen.)

Beachte, dass \textbf{Gemeinfreiheit (\emph{public domain}) keine
gültige Lizenz ist}, weil sie in verschiedenen Ländern unterschiedlich
definiert ist.

Von der WTFPL wird dringend abgeraten\footnote{\url{https://content.luanti.org/help/wtfpl/}}
(englisch); manche verzichten auf deinen Mod, wenn er unter dieser
Lizenz steht.

\subsection{LGPL und CC BY-SA}\label{lgpl-und-cc-by-sa}

Diese Kombination ist in der Luanti-Community verbreitet; auch Luanti
und Minetest Game verwenden sie.

Du stellst deinen Code unter die LGPL 2.1 und deine Grafiken unter CC
BY-SA.

Das bedeutet:

\begin{itemize}
\tightlist
\item
  Jeder darf veränderte und unveränderte Fassungen bearbeiten,
  weitergeben und verkaufen.
\item
  Wer deinen Mod verändert, muss seine Fassung unter dieselbe Lizenz
  stellen.
\item
  Dein Urheberrechtsvermerk muss erhalten bleiben.
\end{itemize}

\subsection{CC0}\label{cc0}

Diese Lizenz eignet sich für Code und Grafiken und erlaubt jedem, mit
deinem Werk zu tun, was er will. Andere dürfen es also verändern,
weitergeben, verkaufen und auch ohne Namensnennung verwenden.

\subsection{MIT}\label{mit}

Eine verbreitete Lizenz für Code. Die einzige Bedingung für Nutzer
deines Codes ist, dass sie in allen Kopien des Codes oder wesentlicher
Teile davon denselben Urheberrechtsvermerk und dieselbe Lizenz beifügen.

\section{Verpacken}\label{verpacken}

Einige Dateien solltest du deinem Mod oder Spiel vor der
Veröffentlichung beilegen.

\subsection{README.txt}\label{readme.txt}

Die README sollte angeben:

\begin{itemize}
\tightlist
\item
  was der Mod oder das Spiel tut und wie man ihn benutzt,
\item
  unter welcher Lizenz er steht,
\item
  optional:

  \begin{itemize}
  \tightlist
  \item
    wo man Probleme melden oder Hilfe bekommen kann,
  \item
    Danksagungen.
  \end{itemize}
\end{itemize}

\subsection{mod.conf / game.conf}\label{mod.conf-game.conf}

Füge unbedingt einen Schlüssel \texttt{description} hinzu, der erklärt,
was dein Mod oder Spiel tut. Fasse dich kurz, ohne vage zu werden. Die
Beschreibung sollte kurz sein, weil sie im Inhaltsinstaller angezeigt
wird, wo wenig Platz ist.

Gutes Beispiel:

\begin{Shaded}
\begin{Highlighting}[]
\NormalTok{description = Fügt Suppen, Kuchen, Gebäck und Säfte hinzu.}
\end{Highlighting}
\end{Shaded}

Vermeide so etwas:

\begin{Shaded}
\begin{Highlighting}[]
\NormalTok{description = Der Essens{-}Mod für Luanti.}
\end{Highlighting}
\end{Shaded}

Das ist zu vage.

\subsection{screenshot.png}\label{screenshot.png}

Screenshots sollten das Seitenverhältnis 3:2 haben (3 Pixel Breite auf 2
Pixel Höhe) und mindestens 300 × 200 Pixel groß sein.

Der Screenshot wird in Luanti als Vorschaubild für den Inhalt angezeigt.

\section{Hochladen}\label{hochladen}

Damit Interessierte deinen Mod herunterladen können, musst du ihn an
einem öffentlich zugänglichen Ort hochladen. Dafür gibt es mehrere Wege.
Nimm den, der für dich am besten passt, solange er diese Anforderungen
erfüllt (und alle weiteren, die Forenmoderatoren vielleicht noch
hinzufügen):

\begin{itemize}
\tightlist
\item
  \textbf{Beständig} -- Die Website sollte nicht ohne Vorwarnung
  abgeschaltet werden.
\item
  \textbf{Direkter Link} -- Ein Klick auf einen Link sollte die Datei
  herunterladen, ohne dass man erst eine andere Seite ansehen muss.
\item
  \textbf{Virenfrei} -- Unseriöse Upload-Dienste enthalten womöglich
  unsichere Werbung.
\end{itemize}

Bei ContentDB\index{contentdb@ContentDB} kannst du ZIP-Dateien
hochladen; ContentDB erfüllt diese Anforderungen.

\subsection{Versionsverwaltung}\label{versionsverwaltung}

Ein Versionsverwaltungssystem (\emph{Version Control System}, VCS)
verwaltet Änderungen an Software und macht es oft leichter, sie zu
verbreiten und Beiträge anderer anzunehmen.

Die meisten Luanti-Modder verwenden Git und eine Plattform wie GitHub,
um ihren Code zu verbreiten.

Git kann anfangs schwierig sein. Hilfe findest du hier:

\begin{itemize}
\tightlist
\item
  Pro Git\footnote{\url{https://git-scm.com/book/de/v2}} -- kostenlos
  online lesbar, auch auf Deutsch.
\end{itemize}

\section{Auf ContentDB
veröffentlichen}\label{auf-contentdb-veruxf6ffentlichen}

ContentDB ist der offizielle Ort, um Inhalte wie Mods, Spiele und
Texturpakete zu finden und zu verbreiten. Benutzer finden Inhalte über
die Website oder laden und installieren sie direkt über das Hauptmenü
von Luanti.

Registriere dich bei ContentDB\footnote{\url{https://content.luanti.org}}
und füge deinen Inhalt hinzu. Lies unbedingt die Hinweise im
Hilfebereich.

\section{Forenthema}\label{forenthema}

Du kannst auch ein Forenthema anlegen, in dem Benutzer über dein Werk
sprechen können.

Themen für Mods gehören ins Forum „WIP Mods``\footnote{\url{https://forum.luanti.org/viewforum.php?f=9}}
(\emph{Work in Progress}), Themen für Spiele ins Forum „WIP
Games``\footnote{\url{https://forum.luanti.org/viewforum.php?f=50}}.
Wenn dein Mod für dich nicht mehr in Arbeit ist, kannst du beantragen,
dass das Thema verschoben wird\footnote{\url{https://forum.luanti.org/viewtopic.php?f=11&t=10418}},
und zwar nach „Mod Releases``.

Das Forenthema sollte ähnliche Inhalte wie die README haben, aber
werbender formuliert sein und einen Link zum Herunterladen enthalten.
Wenn möglich, zeig Screenshots deines Mods in Aktion.

Der Titel des Themas muss eines dieser Formate haben:

\begin{itemize}
\tightlist
\item
  {[}Mod{]} Titel des Mods {[}modname{]}
\item
  {[}Mod{]} Titel des Mods {[}Versionsnummer{]} {[}modname{]}
\end{itemize}

Zum Beispiel:

\begin{itemize}
\tightlist
\item
  {[}Mod{]} More Blox {[}0.1{]} {[}moreblox{]}
\end{itemize}
